\documentclass[a4 paper, 11pt]{article}

\usepackage[ngerman]{babel}
\usepackage[T1]{fontenc}
\usepackage[utf8]{inputenc}

\usepackage{microtype}
\usepackage[top=3cm, bottom=2cm, inner=2cm, outer=2cm]{geometry}

\usepackage{parskip}
\usepackage{amsmath} 
\usepackage{amssymb} 
\usepackage{amsfonts}
\usepackage[locale=DE, per-mode=fraction, quotient-mode=fraction, separate-uncertainty = true]{siunitx}
\usepackage{nicefrac}
\usepackage{wasysym} %Durchmesserzeichen

\usepackage{textcomp}

\usepackage{placeins}

\usepackage{booktabs}
\usepackage{graphics}
\usepackage[rel]{overpic}

\usepackage{caption}

\usepackage[colorlinks, allcolors=black, citecolor=black!80]{hyperref}

\usepackage[all, defaultlines=3]{nowidow}

\usepackage{enumitem}

\usepackage{eurosym}

%\usepackage{comment} % Dient dazu, eine frühere Version bei Bedarf ein- oder ausblenden zu lassen.
%\includecomment{versiona}
%\excludecomment{versiona} 

%\usepackage{color} %farbige Schrift

%\definecolor{new}{RGB}{50, 150, 50} %green
%\definecolor{new}{RGB}{0, 0, 0} %black
%\definecolor{box}{RGB}{50, 150, 50} %green
%\definecolor{old}{RGB}{255, 69, 0} %orange
%\definecolor{move}{RGB}{50,100,255} %blue

%\usepackage{showframe}
%\usepackage{lipsum}

\usepackage{fancyhdr}

\fancyhead[L]{\textsc{Anhang} zum Leitfaden der Maultaschen Tübingen}
%\fancyhead[C]{}
\fancyhead[R]{Version vom \today}

%\renewcommand{\d}{\mathrm{d}}
%\renewcommand{\thesubsubsection}{\alph{subsubsection})}

\title{{\normalsize Ultimate Frisbee in Tübingen} \\ \textbf{\textsc{Anhang} zum Leitfaden der Maultaschen Tübingen}}
\author{Version vom \today \\ {\footnotesize zusammengetippt von Sonja Lorenz\footnote{\href{mailto:so.lorenz@web.de}{so.lorenz@web.de}}}}

\date{}


%TODO: Stehen wirklich nur Sachen drin, die nicht beschlossen werden müssen?

\begin{document}
	
\pagestyle{fancy} % Damit wird Kopfzeile gedruckt.
	
\maketitle

\tableofcontents

\newpage

\section{Turnierfahrten}

\subsection{Vorbereitungen}
Zu Beginn der Planungen schickt der Turnierpate eine Mail mit Teilnehmerliste an den Finanzvorstand (\href{mailto:finanzen@ultimate-tuebingen.de}{finanzen@ultimate-tuebingen.de}) und bekommt von diesem den Link zum Turnierorgasheet zurück, den er an alle Teilnehmer weitergibt. Dort können auf dem ersten Reiter die Fahrtplanung gemacht, Aufgaben verteilt und wichtige Infos/Links weitergegeben werden. 

Bei offiziellen Turnieren kümmern sich Spielervorstand und Turnierpate zusätzlich rechtzeitig (spätestens 3-4 Wochen vorher) darum, dass alle Teilnehmer spielberechtigt und korrekt angemeldet sind (der Typ, der das beim Landesverband macht, kann auch mal im Urlaub sein).

\subsection{Abrechnung und Überweisungssystem}
Der Finanzvorstand spricht:
\begin{quote}
	Idealerweise hat mir der/die Turnierpate/-patin vor dem Turnier bereits eine Mail / Whatsapp / iMessage mit den Teilnehmern des Turniers geschickt und von mir ein Template bekommen, das ihr bereits zur Fahrtplanung etc. genutzt habt. Sobald der/die Abrechnungsbeauftragte dort nach dem Turnier auf dem zweiten Reiter alle Kosten und Fahrten eingetragen hat, sagt diese/r mir Bescheid, sodass ich nochmal schnell drauf schaue. Fertig ist die Abrechnung!
	
	Nochmal in Checklistenform:
	\begin{enumerate}
		\item Mir eine Liste der Teilnehmer schicken (wenn nicht bereits durch Turnierpate erledigt)
		\item Ihr bekommt das Abrechnungssheet (wenn nicht bereits geschehen)
		\item Abrechnungssheet ausfüllen
		\item Ich kontrolliere die Abrechnung 
		\item Offene Beträge können überwiesen werden
	\end{enumerate}
\end{quote}

Außerdem funktioniert die Orga der Überweisungen wie folgt:
\begin{quote}
	JEDER hat ein eigenes Googlesheet mit allen relevanten Infos, die da wären:
	\begin{itemize}
		\item Eigener aktueller Kontostand, negative Beträge sollten möglichst direkt überwiesen werden. Es ist möglich sich einen positiven Kontostand nicht auszahlen zu lassen, um unnötig viele Überweisungen zu vermeiden. Dann wird dieses Guthaben mit den nächsten Ausgaben verrechnet.
		\item Teamkontodaten
		\item alle eigenen Überweisungen
		\item alle gespielten und geplanten Turniere mit Link zum Turnierorgasheet sowie den Kosten/Guthaben daraus
		\item meine E-Mail Adresse (z.B. falls ihr Geld wollt)
	\end{itemize}
	Damit ihr Zugriff auf dieses Sheet bekommt müsst ihr nur eine Sache machen, mir eine E-Mail / Nachricht schreiben (im Idealfall schreibt ihr euren vollen Namen da mit rein, damit ich mit den Namensdopplungen nicht durcheinander komme, außerdem ist eine E-mail als Kontakt am einfachsten zu handhaben).
\end{quote}

\emph{Anmerkung}: Die Emailadresse des Finanzvorstands lautet \href{mailto:finanzen@ultimate-tuebingen.de}{finanzen@ultimate-tuebingen.de} und wäre auf der Maultaschen-Homepage zu finden: \url{https://ultimate-tuebingen.de/}

\subsection{Turnierberichte}
Auf der Maultaschen-Homepage (\url{https://ultimate-tuebingen.de/}) sowie der Facebook-Seite (\url{https://de-de.facebook.com/UltimateTuebingen/}) soll zu jedem Turnier auch ein Bericht veröffentlicht werden. Der Turnierberichtschreiber hat daher die Aufgabe, auf dem Turnier ein Teamfoto zu schießen und im Anschluss möglichst bald einen passenden Text an die entsprechenden Beauftragten (siehe Homepage) zu schicken. Falls der ausführliche Bericht nicht sofort in den nächsten Tagen fertiggestellt werden kann, wäre es gut, zumindest das Foto und einen kurzen Vierzeiler schon mal loszuschicken und den ausführlichen Bericht dann baldmöglichst nachzuliefern.


\section{Bemerkungen zum Trainingsbetrieb}
Das Training für Anfänger wird vollständig über den Hochschulsport angeboten und ist damit der wichtigste \glqq Werbeort\grqq, um neue Mitglieder für das Team zu gewinnen. Daher ist es wünschenswert, dort neben den Grundlagen zum Ultimate auch das soziale Miteinander zu pflegen.

Im nächsten Schritt sind im Verein zunächst das Open- und Damentraining und natürlich das freie Spiel für alle geöffnet, um auf dem Anfängertraining aufbauend weiter an persönlichen Ultimate-Skills und Spieltaktiken zu arbeiten.

Das Mixed-Training ist darüber hinaus grundsätzlich als geschlossene Trainingsgruppe aktiv, in der auf fortgeschrittenem Leistungsniveau, auch mit Blick auf die Vorbereitung der offiziellen Turniere hin, gemeinsam trainiert wird. Dabei liegt es im Ermessen des aktuellen Mixed-Trainers, ob jemand bereits fortgeschritten genug ist, um der Trainingsgruppe beizutreten. 


\section{Eigene Turniere}
Es werden von den Maultaschen aktuell zumeist zwei Turniere im Jahr ausgerichtet, ein Indoor-Turnier im Frühjahr und ein Outdoor-Turnier im Sommer. Diese sind dafür gedacht, interessierte Anfänger in lockerer Atmosphäre ganz unverbindlich ihre ersten Turniere spielen zu lassen, die Teamkasse ein bisschen aufzufüllen und allen ultimativen Spielspaß auf heimischen Boden zu ermöglichen.

Während der Turniere hat das gesamte Team die Rolle des Gastgebers inne und jeder Einzelne ist angehalten, sich auch entsprechend zu verhalten, mitzuorganisieren, mitzuhelfen, Schichten zu übernehmen und die Kuchen- und Salatbar zu bestücken. Auch die Anfänger sollen in die Gastgeberrolle mit eingebunden werden, jedoch in geringerem Umfang (z.B. $1\times$ Salat oder Kuchen und eine Schicht plus Abbau), damit sie genügend Zeit haben, die Turnieratmosphäre auf sich wirken zu lassen.

\subsection{Der Maultaschen-Cup (MTC)}
Der Maultaschen-Cup findet traditionell im Sommer auf dem Gelände des SV03-Stadions bzw. dem Kunstrasenplatz in der Jahnallee statt. Die Teilnehmer dürfen dabei auf dem Stadiongelände zelten und haben während der beiden Turniertage freien Zutritt zum städtischen Freibad direkt gegenüber.

\subsection{Das Live and Let Dive (LALD)}
Das Live and Let Dive ist ein (bisher zumeist eintägiges) Indoor-Turnier im März, das in der Regel über den Hochschulsport organisiert wird und in einem solchen Fall auch nur für Hochschulteams geöffnet ist. Es findet häufig in der Unisporthalle in der Alberstraße, Tübingen statt und ist, unter anderem durch das Format des Ein-Tages-Turniers, besonders gut dafür geeignet, (die eigenen) Anfänger erste Turnierluft schnuppern zu lassen.


\section{Bemerkungen zur Trikotgestaltung}
Seit 2009 sind die Farben der Maultaschen und die Trikotfarben der Mixed-Mannschaft gelb und lila. Jeweils eine davon entspricht (in Akzenten) auch den Trikotfarben der Herren (gelb) und der Damen (lila).
 
Um die Zusammengehörigkeit der drei Mannschaften (\glqq Maultaschen Tübingen\grqq, \glqq THW\grqq\ und \glqq MissTürious\grqq) und die Zugehörigkeit zum Verein noch weiter zu unterstreichen, wäre es erstrebenswert, in Zukunft auch ein gemeinsames Logo sowie das Logo des Vereins auf allen Trikots abzubilden.

Mit Blick auf die Teilnahme an offiziellen Turnieren ist es zusätzlich empfehlenswert, die Trikots gemäß der aktuellen WFDF/dfv-Richtlinien zu entwerfen (helles und dunkles Trikot, Platzierung der Nummer,~\dots).


\section{Geschenke}
Es ist zur Tradition geworden, dass alle Teammitglieder, die die Maultaschen leider verlassen müssen, zum Abschied ein Maultaschen-Handtuch überreicht bekommen. Außerdem erhalten alle frischgebackenen Eltern eine Maultaschen-Babyhose für den Nachwuchs. Diese beiden Dinge werden aus der gemeinsamen Teamkasse finanziert. 


Es gibt auch Geschenke/Aufmerksamkeiten für besonderes Engagement und besonderen Einsatz im Team bzw. im Team-/Vereinsumfeld. Dabei geht es uns vor allem um Anerkennung, Wertschätzung und Dankbarkeit, weshalb wir hier einen Rahmen schaffen wollen, wie solche Geschenke gehandhabt werden. Beispielhafte Gründe für Geschenke sind: Turnierorga, sonstige Orga von größeren Teamevents (Winter-Wochenende), Arbeit an der Vereinsinfrastruktur (Leitfaden, Homepage usw.). 

Geschenke bis 10\euro, an Personen innerhalb des Team-/Vereinslebens, bedürfen eines speziellen Grundes (s.o.). Zudem muss die Idee des Geschenks von mind. 3 Vereinsmitglieder einschließlich mind. 1 Vorstandsmitglied vor dem Kauf unterstützt werden. Vorstände haben die Möglichkeit den Kauf aufzuschieben, um den gesamten Vorstand zu befragen. \\
Bei Geschenken von 10\euro-30\euro bedarf es der Zustimmung des gesamten Vorstandes und zusätzlich mind. 7 Vereinsmitglieder, die das Geschenk befürworten.\\ 
Über 30\euro sollten Geschenke nur selbständig organisiert und durch Einsammeln finanziert werden. 

Die Kosten der Geschenke werden aus der gemeinsamen Teamkasse gezahlt.

Die Übergabe von Geschenken sollte grundsätzlich an einem Teamevent (Training, Sitzung, …) erfolgen. Außerdem wird bei der Teamsitzung im Februar der ausgegebene Gesamtbetrag und die Anzahl der Geschenke vom Vorjahr bekannt gegeben. Der Finanzvorstand führt eine detaillierte Liste der Geschenke, die auf Anfrage bei ihm eingesehen werden kann.

\section{Homepage und aktuelle Infos}
Alle aktuellen Infos (Trainingszeiten \& -orte, aktuelle Vorstände und weitere verantwortliche Personen,~\dots), sowie nützliche Unterlagen (Dokus zu den ausgerichteten Turnieren, Protokolle, Formular zum Fahrtkostenzuschussantrag bei der Stadt, der Leitfaden \& dieser Anhang,~\dots) sind (teilweise nur auf dem internen Teil) der Homepage unter \url{https://ultimate-tuebingen.de/} zu finden und werden vom Homepage-Beauftragten auf dem neuesten Stand gehalten. 

Die Online-Abstimmungen bei weitreichenden Entscheidungen oder  für die Trainerlegitimierung finden ebenfalls über den internen Bereich der Homepage statt.

Dazu spricht der Homepage-Beauftragte:
\begin{quote}
Es gibt jetzt einen internen Bereich
für Spieler mit verschiedenen Inhalten und Funktionen. Turniere können
von der Turnier-Orga eingetragen werden und Spieler können sich nach dem
gewohnten Dabei-Fraglich-Nichtdabei-Prinzip anmelden.

Dafür muss sich jeder registrieren. Den Link zum Einloggen/Registrieren
findet ihr ganz unten am Ende der Website oder direkt hier:

Login: \url{https://ultimate-tuebingen.de/wp-login.php}

Registrieren: \url{https://ultimate-tuebingen.de/wp-login.php?action=register}
(haltet ein schönes Profilbild im quadratischen Format bereit)

Das ist jetzt gar nicht ausschließlich für die aktivsten
Turnierteilnehmer interessant, sondern für alle, die aktiv im Verein
bzw. beim Ultimate in Tübingen einbringen wollen. Ihr werdet anschließend
manuell freigeschaltet, jedoch müsst ihr mir hierfür nicht Bescheid geben.
\end{quote}


\section{Bemerkung zur Stimmberechtigung}
Das im Leitfaden beschlossene Stimmrecht ist bewusst auf Vereinsmitglieder beschränkt, um ein Kriterium zu haben, das eindeutig festgelegt und für alle Entscheidungen/Wahlen einheitlich anwendbar ist. Da die Orga des Hochschultrainings aus dem Verein heraus passiert, und beide Dinge zu sehr verwoben sind, um sie sauber zu trennen (der Unisport würde ohne den Verein in dieser Form wohl nicht existieren), sollte es vertretbar sein, die Mitwirkung engagierter Unisportler, die mitgestalten, aber nicht in den Verein eintreten wollen, auf die Möglichkeiten zu beschränken, die im Rahmen einer aktiven Teilnahme an den Teamsitzungen gegeben sind.


\section{Verantwortlichkeiten der einzelnen Vorstandsposten} \label{ag_vorstand}

\subsection{Abteilungsvorstand}
Kommunikation mit Verein:
\begin{itemize}
	\item Teilnahme an den Abteilungsleitersitzungen (immer am 1. Dienstag im Monat von 19:30~Uhr bis ca. 21~Uhr)
	\item verpflichtende Teilnahme an der ordentlichen Mitgliederversammlung des SV03 (alle 2 Jahre im Juni)
	\item Annahme von Anfragen zu Mitgliedern, Abrechnungen, Hilfsdiensten (Erbelauf, Tigers
	Dance Night, \dots)
	\item Entgegennahme von Abrechnungen vonseiten des Vereins
	\item Dokumentation und Kontrolle der Mitglieder zusammen mit dem Spielervorstand (Mitgliederliste wird auf Anfrage vom Verein zugesendet)
	\item Vertreten der Belange des Frisbeeteams im Verein
	\item Teilnahme an Workshops des Vereins
\end{itemize}
\vspace{0.3cm}
Vereins-Frisbeeturniere initiieren (MTC, DM-Ausrichtungen, \dots):
\begin{itemize}
	\item Terminabsprache mit Fußball (bei Rasenturnier)
	\item Buchung der Plätze (Jahnallee und Nebenplatz) oder Halle (z.B. Pfullingen) (teils
	zusammen mit Trainervorstand, der Kontakt zur Stadt hat)
	\item Orgateam finden und für Fragen zur Verfügung stehen
	\item ggf. Unterstützung bei der Organisation der Infrastruktur bei Peter (Kreide, Biertischgarnituren, Jugendraum, Küche, \dots)
	\item Bezahlung der Infrastruktur an Fußballer (Umbuchung vom Abteilungskonto)
	\item Sicherstellung, dass die Abrechnung pünktlich beim Verein eingeht
\end{itemize}

\subsection{Hochschulsportvorstand}
Kommunikation mit Hochschulsport:
\begin{itemize}
	\item Absprache der Trainingszeiten und Hallen- bzw. Platzvergabe
	\item Ansprechpartner für Ingrid Arzberger [Wer ist das?]
	\item Ansprechpartner für Studenten bzgl. Trainingsbetrieb
	\item Beantwortung von Anfragen von Extern (Schulen, Vereine, \dots)
\end{itemize}
\vspace{0.3cm}
Organisation der Teilnahme an der deutschen Hochschulmeisterschaft (DHM):
\begin{itemize}
	\item rechtzeitige Anmeldung (noch vor Semesterbeginn)
	\item ggf. den Zusammenschluss von Spielgemeinschaften initiieren (muss mit noch mehr Vorlauf passieren)
	\item Kader auswählen bzw. verantwortliche Personen benennen
\end{itemize}
\vspace{0.3cm}
Hochschulsport-Frisbeeturniere initiieren (LALD und andere, auch DHM):
\begin{itemize}
	\item Terminvergabe und Hallenreservierung
	\item Orgateam finden und unterstützen
	\item Kontakt zu Hausmeister etc.
\end{itemize}
\vspace{0.3cm}
Halbjährliche große Teamsitzung organisieren (im Herbst und im Frühjahr zu Semesterbeginn):
\begin{itemize}
	\item Terminplanung
	\item Raumplanung
	\item Ankündigung an das Team
	\item Aufstellung der Agenda
	\item im Herbst: Sicherstellung, dass es ausreichend viele Bewerber für die Vorstandsposten gibt
	\item Durchführung (ggf. Sitzungsleiter ernennen)
	\item Protokoll (Protokollanten beauftragen)
\end{itemize}

\subsection{Finanzvorstand}
\begin{itemize}
	\item Kontrolle der Finanzen der Abteilung (Bezahlung Plätze, Flutlicht, \dots)
	\item Bezahlung der Kabinennutzung an die Fußballabteilung
	\item Abgabe der Turnierabrechnungen beim Verein (wenn sie über den Verein ausgerichtet wurden)
	\item Kontrolle der Turnierfahrtenabrechnungen
	\item Verwaltung der Einzahlungen und Auszahlung von Teammitgliedern
	\item (Aus-)Zahlung von Sonderausgaben (Trikots, Scheiben, \dots)
	\item Kontrolle der Reisekostenumverteilung des dfv bzw. Zahlung der Rückerstattung
	\item Kontrolle der Fahrtkostenzuschüsse der Stadt
	\item Erstellung eines Jahresfinanzberichts
	\item Finden und Koordination von Kassenprüfern (Legitimation der Kassenprüfer in der Teamsitzung im Winter, Bericht zur Kassenprüfung im April)
\end{itemize}

\subsection{Trainervorstand}
Kommunikation mit Stadt:
\begin{itemize}
	\item Ansprechpartner für Frau Dettmer (Platzvergabe)
	\item Organisation von Schlüsseln, Marken
	\item Organisation von Bezahlungen in Absprache mit dem Finanzvorstand (Plätze, \dots)
	\item Annahme von Anfragen (Tauschen von Terminen, Rasenrenovation, Probleme auf
	Jahnallee, \dots)
\end{itemize}
\vspace{0.3cm}
Organisation des Trainingsbetriebs:
\begin{itemize}
	\item Absprache mit Fußball, wann und wie Training möglich ist
	\item Buchung aller Plätze für Training (Jahnallee, Stadion)
	\item Buchung des Platzes beim SSC inklusive Bezahlung
	\item ggf. Buchung weiterer Plätze/Hallen
	\item alle Trainer koordinieren, sodass taktisch und philosophisch an einem Strang
	gezogen wird
	\item Trainer für die verschiedenen Trainings finden
\end{itemize}
\vspace{0.3cm}
Halbjährliche kleine Teamsitzung organisieren (im Winter und im Sommer):
\begin{itemize}
	\item Terminplanung
	\item Raumplanung
	\item Ankündigung an das Team
	\item Aufstellung der Agenda
	\item Durchführung (ggf. Sitzungsleiter ernennen)
	\item Protokoll (Protokollanten beauftragen)
\end{itemize}

\subsection{Spielervorstand}
Offizieller Spielbetrieb:
\begin{itemize}
	\item offizieller dfv-Kontakt (Onlinetool)
	\item Meldung der Spieler beim dfv und Sicherstellung, dass sie in einem Verein sind und
	die Datenschutzerklärung bzw. Athletenvereinbarung vorliegt
	\item Initiator der Bezahlung der Saisongebühren (Kommunikation mit dem Finanzvorstand)
	\item Anmeldung der Teams auf dfv-turniere.de (unter Turniere)
	\item Unterstützung der Turnierpaten der offiziellen Turniere (Spielberechtigungen sicherstellen
	durch DSE und AV, allgemeine Orga, Spielerlisten pünktlich einstellen (auf dfv-turniere.de), \dots)
	\item Annahme von Anfragen des dfv (Saisongebühren, Turnierausrichtungen, \dots)
	\item Anmeldungen für Verein einsammeln und an Abteilungsvorstand oder direkt an den Verein übergeben
\end{itemize}
\vspace{0.3cm}
Koordination weiterer Jobs im Frisbeeteam:
\begin{itemize}
	\item Wissen, wer aufhört und neuen Verantwortlichen finden als:
	\begin{itemize}
		\item Turnierkoordinatoren (via Onlinetool)
		\item Homepage-Beauftragter (Vereinsseite und eigene Homepage)
		\item Facebook-Beauftragter
		\item Lichtmarkenverwalter
		\item Markenbring- und Schließbeauftragter fürs freie Spiel
		\item Scheibenverkäufer
		\item Pickup-Trikot-Verwalter
		\item Leitfaden-Verantwortlicher
		\item Maultaschen-Boxen-Mami
		\item ggf. Scheibendesign-Team
		\item ggf. Trikot/Lifestyleware-Team
	\end{itemize}
\end{itemize}
\vspace{0.3cm}
Diverses:
\begin{itemize}
	\item Abschieds-Handtücher überreichen und wenn nötig nachbestellen (sprich: Überblick
	behalten wer geht)
	\item Beantwortung von Anfragen über Emailadresse der Homepage
	\item Hütchen und Leibchen nachbestellen
	\item Socialevents planen und koordinieren (z.B. Winterfest)
	\item Ansprechpartner für freies Spiel (Emailanfragen beantworten, Ansprechpartner
	für Änderungswünsche am Spiel-System)
\end{itemize}


\section{Änderungen an diesem Anhang zum Leitfaden}
\dots können vom Leitfadenbeauftragten jederzeit durchgeführt werden, wenn ihm Ergänzungen oder ein inhaltliches Update (z.B. nach neuer Diskussion eines Themas in einer Teamsitzung) sinnvoll erscheinen.

Die alte Version wird dann einfach online (im internen Bereich der Homepage) durch die aktualisierte ersetzt.

\end{document}

